
\section{Related Work}
\subsection{Audio-driven Character Animation}
Pioneering audio-driven character animation utilized multi-stage pipelines with intermediate representations like landmarks, depth maps, and neural fields \cite{zhongIdentityPreservingTalkingFace2023, wilesX2FaceNetworkControlling2018, hongDepthAwareGenerativeAdversarial2022, guoADNeRFAudioDriven2021, zhouyangMakeltTalk2020, prajwalLipSyncExpert2020}. 
These intermediate steps simplified learning by modeling structure explicitly, but introduced error accumulation and reduced visual realism.
Recently, diffusion models make direct synthesis feasible by implicitly learning geometry, identity, and motion, thus removing the need for intermediate steps and enhancing realism. \cite{wangFantasyTalkingRealisticTalking2025, weiMoChaMovieGradeTalking2025, yiMagicInfiniteGeneratingInfinite2025, cuiHallo3HighlyDynamic2025, chenEchoMimicLifelikeAudioDriven2024, yangMegActor$S$UnlockingFlexible2024}. Although achieving impressive results, these methods rely on pre-defined inpainting frames—typically the first frame of the video used to guide generation—limiting scene flexibility. Besides, they focus primarily on single-character scenarios restricting multi-character applications, as shown in Table \ref{tab:comparison-of-supported}. Concurrent works \cite{kong2025let,chenHunyuanVideoAvatarHighFidelityAudioDriven2025} attempt to extend talking head generation to multi-character scenarios, but they still rely on a provided inpainting frame. In addition, their use of static masks limits the modeling of dynamic scenes.
\subsection{Multi-Concept Video Generation}
Recent advances enable multi-concept video generation through specialized architectures. MAGREF \cite{dengMAGREFMaskedGuidance2025} introduces masked guidance for multi-subject synthesis, while DanceTogether \cite{chenDanceTogetherIdentityPreservingMultiPerson2025} preserves identities in interactive scenes. Concat-ID \cite{zhongConcatIDUniversalIdentityPreserving2025} and ConceptMaster \cite{huangConceptMasterMultiConceptVideo2025} achieve universal identity preservation using concatenated embeddings and decoupled concept learning.
Ingredients \cite{fei_ingredients_2025} project facial embeddings and visual tokens into a shared space to compute relevance scores for assigning each token to a character. However, the resulting masks lack precision and temporal smoothness due to the absence of geometric priors that naturally constrain 3D mask shapes.
Moreover, these methods primarily rely on image and text conditioning and lack support for audio-driven multi-character generation. This limitation restricts their applicability in scenarios that require audio-synchronized narrative content creation.
\begin{figure*}[t]
    \centering
    \includegraphics[width=1.01\textwidth]{figure2} % Reduce the figure size so that it is slightly narrower than the column. Don't use precise values for figure width.This setup will avoid overfull boxes.
    \caption{Illustration of the three approaches to implement the \textbf{Embedding Router}: (1) \textbf{Pre-Denoise Router}, which uses a static visual mask extracted from the inpainting frame; (2) \textbf{Post-Denoise Router}, which detects face regions after the initial denoising process; and (3) \textbf{Intra-Denoise Router}, which dynamically generates visual masks during denoising using a learned module.}
    \label{fig2}
\end{figure*}
